\section{Evaluation protocol and source audit}
\label{app:protocol}

This appendix gives the detailed protocol behind the primary comparison.
The main text uses only DiffusionGemma-26B-A4B and Qwen3-30B-A3B;
smaller models appear only in explicitly secondary experiments.

\paragraph{Models and serving.}

Both primary models use a mixture-of-experts (MoE) design, activating a
subset of their expert parameters for each token. Their broadly comparable
parameter and weight footprints control architecture class and approximate
scale, but do not match training data, model quality, tokenizer, or decoder
implementation. The main measurements use the pinned vLLM 0.25.1 stack.
Earlier Qwen2.5 experiments use smaller dense autoregressive models,
which do not select expert subsets, including Qwen2.5-3B (dense AR-3B)
and Qwen2.5-1.5B (dense AR-1.5B). LLaDA-8B is a separate diffusion
baseline. These secondary comparisons are identified throughout.

\paragraph{Task and data.}
Each raw log line is mapped to a JavaScript Object Notation (JSON) record,
a structured set of named fields and values, with twelve fields. Seven
structural fields describe the timestamp, host, component, severity,
message template, error code, and numeric value; two fields encode an
anomaly label and score; three provide semantic or action information.
The primary corpus is BGL from the Loghub collection~\cite{loghub}, containing
4{,}747{,}963 BlueGene/L log lines. The preparation procedure samples
600{,}000 lines by reservoir sampling (seed 42) and splits the sample
temporally into 400{,}000 training and 200{,}000 test records.
Evaluation cells use subsets of this test split, usually $n{=}300$ unless
otherwise stated. Hadoop Distributed File System (HDFS) logs provide a
secondary log-corpus check using the
2{,}000-line Loghub sample.

Reference records are constructed by the repository's parsing pipeline,
which supports the Drain log-template parser~\cite{drain} and a
regular-expression fallback. They are not
independently annotated semantic ground truth. In particular, the code
sets \texttt{is\_anomaly} when the native alert tag is non-normal
\emph{or} severity is ERROR, FATAL, or CRITICAL, and sets
\texttt{anomaly\_score} to the corresponding binary value. Thus these
fields measure agreement with a label-construction rule, not solely
prediction of native alert tags. A FATAL line tagged \texttt{-}, for
example, still receives an anomaly label under this rule.

\paragraph{Metrics and decoding.}
Quality is macro-averaged token-overlap F1 across seven structural and two
anomaly fields. We exclude \texttt{event\_type},
\texttt{affected\_entity}, and \texttt{remediation\_hint}, whose
references are parser-derived rather than independently validated.
Excluding them reduces one source of contamination; it does not make the
remaining labels independent of the preparation pipeline.
F1 is the harmonic mean of precision and recall; here these measure token
overlap between predicted and reference field values. Macro-averaging gives
the scored fields equal weight. The scorer treats two empty field values
as a match, so Field-F1 should be
read alongside validity. \emph{Schema validity} is the harness check for
required fields and supported type and allowed-value constraints on the parsed record;
\emph{completeness} measures field presence. Neither establishes semantic
correctness. We use first-object-wins parsing for repeated JSON outputs.

Both models use greedy decoding, selecting the highest-scoring token
without random sampling (temperature 0, top-p 1, repetition penalty 1).
Unless varied explicitly, DG uses $K{=}8$ denoising steps, the successive
refinement passes within a canvas, and canvas size $C{=}256$ token positions.
Model-specific settings and environment versions
are recorded in the result artifacts. Nominal model sizes are 26B total/4B
active parameters for DG and 30B total/3B active for AR, with approximate
bfloat16 (bf16, a 16-bit floating-point format) weight footprints of 52 and
60\,GB; these are rounded model labels,
not measurements of total runtime memory.

We distinguish three quantities: generated output length $L$, the configured
maximum output budget $B$, and the number of decoded canvases
$N_C$. Raising $B$ does not increase $L$ when a model has already emitted
end-of-sequence (EOS). We also distinguish serial mean service cost,
per-event latency percentiles, and inverse throughput under batching.
These quantities answer different deployment questions. The $p50$ is the
median; $p99$ is the latency at or below which 99\% of observations fall.
A confidence interval (CI) describes statistical uncertainty in an estimate.
The coefficient of variation (CV) is standard deviation divided by mean,
reported as a percentage here.


\subsection{Paired quality reanalysis}
We re-score the first 300 saved predictions from the DG repeat-1 file
against the AR-30B seed-42 predictions. Both runners read the first
records of the same prepared BGL test split without shuffling at seed 42.
The other AR seeds permute that same 300-record pool; they are repeated
timing runs, not 900 independent quality examples.

A frozen copy of the 300 parser-derived reference records accompanies this
snapshot. The reanalysis uses the repository's nine-field scorer and schema
checker without changes. It reproduces AR-30B's archived F1 to its recorded
precision. The DG subset is drawn from the 2,000-event run; it is not a new
model run or the separate 300-event budget sweep. No quality-based filtering
selects records.

For each event we subtract AR's nine-field F1 from DG's. A paired percentile
bootstrap resamples these 300 differences with replacement 10,000 times
using NumPy's default generator at seed 42. The mean difference is
$-0.0181$, with a 95\% interval of $[-0.0289,-0.0086]$.
This conditional, post-hoc analysis shows a small DG deficit on this metric.
No equivalence margin is selected after observing the result. Shared log
templates and parser-derived references limit generalization.

\texttt{audit\_primary\_pair.py} produces the quality and budget tables
from saved outputs. Its JSON report records source hashes, per-field scores,
sample sizes, and the bootstrap settings. The reported schema-valid fraction
checks the harness's supported constraints, not arbitrary JSON-schema
compliance or semantic correctness.

\subsection{Software provenance}
The intended stack is vLLM 0.25.1, torch 2.11.0, CUDA 13.0, on an A100 80GB.
The AR-30B summaries record that stack, as does
\texttt{results/anchors/anchor\_remeasure.json} for DG.
However, \texttt{results/anchors/dg\_L256\_n2000.json} reports the same
DG timing summary with vLLM 0.11.0, torch 2.8.0, and CUDA 12.8.
The saved records do not resolve this metadata conflict. We retain both
source artifacts and describe the timing result as a comparison of saved
runs, not proof of exact software parity.

\subsection{Timing uncertainty}

In the repeated BGL runs at batch size one, AR-30B has median latency
\arMedian\,ms (95\% bootstrap CI [\arCiLo, \arCiHi]; 900 events
pooled across seeds 42, 137, and 314). DG has median \dgMedian\,ms
(CI [\dgCiLo, \dgCiHi]; 2,000 events). The ratio of these medians
is approximately $2.4\times$. Bootstrap intervals use 10,000 event-level
resamples. The AR seed-median CV is \seedCV\%, describing repeatability
across these three runs rather than uncertainty across hardware or workloads.

A two-sided permutation test of the absolute median difference observed
no equally or more extreme result in 100,000 shuffles. The finite-sample
Monte Carlo estimate $(0+1)/(100000+1)$ is approximately $10^{-5}$;
we do not interpret the stored zero exceedance fraction as a zero
probability. This unpaired test treats event timings as exchangeable and
does not control run-level differences or temporal dependence. The
observed latency gap is large, but its interval is conditional on these
measurement assumptions.


The blue band in an earlier version of the cumulative latency figure
marked a confidence interval for the median. It was not another population
of responses. The revised figure omits these bands and marks completion
fractions at a stated generation-time budget instead.

The budget table is a post-hoc count of saved generation times, including
all recorded events. A completion at exactly the budget is counted as on
time. A second count also requires harness schema validity. Neither count
tests semantic correctness or includes queueing, tools, and downstream work.
